% ===========================================================================
%  离散数学与结构（I） 2026-10-13（第 5--6 节）课堂笔记
%  环、理想与同构定理
%  Compile:  latexmk -xelatex main      (or: xelatex main, twice)
% ===========================================================================
\documentclass[cjk]{loom}

\newcommand{\N}{\mathbb{N}}
\newcommand{\Z}{\mathbb{Z}}
\newcommand{\Q}{\mathbb{Q}}
\newcommand{\R}{\mathbb{R}}
\newcommand{\C}{\mathbb{C}}
\newcommand{\HH}{\mathbb{H}}               % 四元数
\newcommand{\F}{\mathbb{F}}                % 有限域
\newcommand{\ideal}{\trianglelefteq}       % 理想
\newcommand{\Frac}{\mathrm{Frac}}          % 分式域
\newcommand{\charr}{\mathrm{char}}         % 特征（\char 是 TeX 原语，故用此名）
\newcommand{\Img}{\mathrm{Im}\,}           % 像
\newcommand{\End}{\mathrm{End}}            % 自同态环

\runningthread{离散数学与结构 · 第四周（上）}

\begin{document}

\loomcover
  {环、理想与同构定理}
  {离散数学与结构（I）· 第四周（上）课堂笔记}
  {}
  {2026 年 10 月 13 日}

\begin{strand}
本讲脉络：先给出环的定义与术语（含幺环、交换环、除环、域、零因子、整环）；
然后是一批环的例子——数环、$\Z/n\Z$、矩阵环、自同态环、函数环、多项式环、
群环、四元数；接着讨论子环、环同态、核与理想、商环；证明第一同构定理并
建立理想的运算（交、和、积、生成理想与主理想）；再证第三同构定理与对应
定理，刻画单环——含幺交换环是域当且仅当它没有非平凡理想；由此得到极大
理想与素理想的刻画（$M$ 极大 $\iff$ $R/M$ 是域，$P$ 素 $\iff$ $R/P$
是整环）；随后是中国剩余定理（环版本）；最后讨论单位与单位群、分式域、
特征，并构造有限域 $\F_p$ 与无穷特征 $p$ 域 $\F_p(x)$。

本讲的动机：抽象代数部分的最终目标是讨论\keyword{有限域}以及性质非常好
的环（特别是多项式环）。
\end{strand}

\clearpage
{\small\tableofcontents}

\clearpage
% ===========================================================================
\section{环的定义}\label{sec:rings}

\subsection{环与含幺环}

\begin{definition}[环（Ring）]
环 $(R, +, \cdot)$ 是一个集合 $R$ 连同两个二元运算——\keyword{加法} $+$
与\keyword{乘法} $\cdot$（常省略记 $ab$），满足：
\begin{enumerate}[leftmargin=1.6em,itemsep=2pt]
  \item $(R, +)$ 是\keyword{阿贝尔群}：结合律、交换律、存在零元 $0$、
        每个 $a$ 有负元 $-a$；
  \item $(R, \cdot)$ 是\keyword{半群}：结合律 $(ab)c = a(bc)$；
  \item \keyword{分配律}（乘法对加法）：
        \[
          a(b+c) = ab + ac, \qquad (a+b)c = ac + bc.
        \]
\end{enumerate}
\end{definition}

\begin{definition}[含幺环]
若存在 $1\in R$ 使对任意 $a\in R$ 有 $1\cdot a = a\cdot 1 = a$，则称 $R$
为\keyword{含幺环}。
\end{definition}

\begin{remark}[术语说明]
环的定义在文献中不统一：一派把不一定含幺的结构叫环（ring），含幺的叫
含幺环；另一派把含幺的叫环，不一定含幺的叫 \keyword{rng}（把 ``ring''
中的 ``i''（identity）丢掉）。后面真正关心的对象（如域）几乎都是含幺
且交换的。
\end{remark}

\begin{example}[无幺元的环]
在任一阿贝尔群上定义乘法 $a\cdot b = 0$（任意两元素相乘恒为零），则除
“存在幺元”外其余公理全部满足。
\end{example}

\subsection{交换环、除环、域}

\begin{definition}[交换环]
乘法满足交换律 $ab = ba$ 的环称为\keyword{交换环}。
\end{definition}

\begin{definition}[除环（Division Ring）]
若含幺环 $R$ 中每个非零元素都有乘法逆，即对任意 $a\ne 0$ 存在 $a^{-1}$
使
\[
  a\cdot a^{-1} = a^{-1}\cdot a = 1,
\]
则称 $R$ 为\keyword{除环}（体）。
\end{definition}

\begin{definition}[域（Field）]
交换的除环称为\keyword{域}。即：含幺、交换、每个非零元可逆。
\end{definition}

\subsection{零因子与整环}

\begin{definition}[零因子（Zero Divisor）]
环 $R$ 中若存在 $a\ne 0$，$b\ne 0$ 使得 $ab = 0$，则称 $a$ 为
\keyword{左零因子}、$b$ 为\keyword{右零因子}，统称\keyword{零因子}。
\end{definition}

除环（特别地，域）中没有零因子：若 $ab = 0$ 且 $a\ne 0$，两边左乘
$a^{-1}$ 即得 $b = 0$。

\begin{definition}[整环（Integral Domain）]
无零因子的交换（含幺）环称为\keyword{整环}。
\end{definition}

% ===========================================================================
\section{环的例子}\label{sec:ring-examples}

\begin{example}[数环]
$\Z$、$\Q$、$\R$、$\C$ 在通常加法与乘法下都是含幺交换环。
\end{example}

\begin{example}[模 $n$ 剩余类环]
$\Z/n\Z$ 是含幺交换环；$n$ 为合数时有零因子。
\end{example}

\begin{example}[矩阵环]
设 $R$ 是环，$R$ 上全体 $n\times n$ 矩阵 $M_n(R)$ 在通常矩阵加法与乘法
下构成环。$R$ 含幺当且仅当 $M_n(R)$ 含幺（单位矩阵）；即使 $R$ 交换，
$M_n(R)$（$n\ge 2$）一般也不交换。
\end{example}

\begin{example}[自同态环]
设 $A$ 为阿贝尔群，全体自同态 $\End(A)$ 构成环：加法逐点定义
$(f+g)(a) = f(a)+g(a)$，乘法为复合 $f\circ g$。
\end{example}

\begin{example}[直积]
环的笛卡尔积 $R_1\times R_2\times R_3$ 在按分量的运算下构成环。
\end{example}

\begin{example}[函数环]
对任意集合 $S$ 与环 $R$，全体映射
\[
  R^S = \{\,f\colon S\to R\,\}
\]
在逐点运算下构成环。取 $S = \{1,2,3\}$ 时 $R^S$ 就是直积
$R\times R\times R$。
\end{example}

\begin{example}[多项式环]
设 $R$ 为环，$X$ 为形式符号，定义\keyword{多项式环}
\[
  R[X] = \left\{ f(X) = c_0 + c_1 X + \cdots + c_n X^n
    \ \middle|\ c_i\in R \right\},
\]
元素是只有有限项非零的形式多项式；加法对应系数相加，乘法按形式展开
合并同类项。一般讨论中要求 $R$ 含幺且交换，但单就写出定义而言这两个
条件都不必要。
\end{example}

\begin{example}[群环 / 半群环]
设 $R$ 为环，$G$ 为半群，定义\keyword{半群环} $R[G]$：元素为形式有限和
\[
  \sum_{g\in G} r_g\, g, \qquad r_g\in R\ (\text{仅有限个非零}),
\]
加法对应项相加，乘法由分配律展开（系数在 $R$ 中相乘、半群元素在 $G$
中相乘）：
\[
  \left(\sum_g r_g g\right)\left(\sum_h s_h h\right)
  = \sum_{g,h} r_g s_h\, (gh).
\]
取半群 $G = \{X^0, X^1, X^2, \dots\}$（$X^i\cdot X^j = X^{i+j}$），
$R[G]$ 正是多项式环 $R[X]$。
\end{example}

\begin{example}[Hamilton 四元数]
取四个符号 $1, i, j, k$，规定
\[
  i^2 = j^2 = k^2 = -1, \qquad
  ij = k,\ ji = -k,\ jk = i,\ kj = -i,\ ki = j,\ ik = -j.
\]
\keyword{实四元数环} $\HH$ 的元素形如 $a + bi + cj + dk$
（$a,b,c,d\in\R$）。$\HH$ 的特点：
\begin{itemize}[leftmargin=1.3em,itemsep=2pt]
  \item \keyword{不交换}：$ij = k$ 而 $ji = -k$；
  \item 但它是\keyword{除环}：$a+bi+cj+dk\ne 0$ 时
        \[
          (a+bi+cj+dk)^{-1}
          = \frac{a - bi - cj - dk}{a^2+b^2+c^2+d^2}.
        \]
\end{itemize}
这是一个“不交换但每个非零元可逆”的除环（因而不是域）的例子。
\end{example}

% ===========================================================================
\section{子环、同态与理想}\label{sec:subring-hom-ideal}

\subsection{子环与环同态}

\begin{definition}[子环]
$R'\subseteq R$ 称为 $R$ 的\keyword{子环}，若它在同样的加法与乘法下
自身构成环。验证时只需检查 $R'$ 对\keyword{减法与乘法封闭}。
\end{definition}

\begin{definition}[环同态]
映射 $f\colon R\to R'$ 称为\keyword{（环）同态}，若对任意 $a,b\in R$：
\[
  f(a+b) = f(a) + f(b), \qquad f(ab) = f(a)f(b).
\]
由于 $f$ 保加法，它自动是加法群同态，故 $f(0) = 0$，$f(-a) = -f(a)$。
\end{definition}

\begin{proposition}[像是子环]
同态 $f\colon R\to R'$ 的像 $\Img f = \{\,f(a) : a\in R\,\}$ 是 $R'$
的子环。
\end{proposition}

\begin{proof}
由群论已知 $\Img f$ 对减法封闭；又 $f(a)f(b) = f(ab)\in\Img f$，
对乘法封闭。
\end{proof}

\subsection{核与理想}

\begin{definition}[核]
环同态 $f\colon R\to R'$ 的\keyword{核}为
\[
  \ker f = \{\,b\in R : f(b) = 0\,\}.
\]
\end{definition}

核是子环，且还有更强的性质：对任意 $a\in R$、$b\in\ker f$，
\[
  f(ab) = f(a)f(b) = f(a)\cdot 0 = 0 \;\Longrightarrow\; ab\in\ker f,
\]
同理 $ba\in\ker f$。即\keyword{环中任意元素乘核中元素（无论左右）仍
落在核中}。

\begin{definition}[理想（Ideal）]
环 $R$ 的非空子集 $I$ 称为 $R$ 的\keyword{理想}，记 $I\ideal R$，若
\begin{enumerate}[leftmargin=1.6em,itemsep=2pt]
  \item $(I,+)$ 是 $(R,+)$ 的子群（对减法封闭）；
  \item \keyword{吸收性}：对任意 $a\in R$、$b\in I$，有 $ab\in I$ 且
        $ba\in I$。
\end{enumerate}
\end{definition}

\begin{strand}
理想的定义正是\keyword{照着核的性质}写出来的，因此同态的核自然是
理想——理想在环论中的地位对应正规子群在群论中的地位。
\end{strand}

\subsection{商环（Quotient Ring）}

设 $I$ 是环 $R$ 的理想。$(I,+)$ 是阿贝尔群 $(R,+)$ 的子群（自动正规），
可考虑加法陪集
\[
  a + I = \{\,a + i : i\in I\,\},
\]
其全体记为 $R/I$。在 $R/I$ 上定义运算：
\[
  (a+I) + (b+I) = (a+b) + I, \qquad (a+I)(b+I) = ab + I.
\]

\begin{theorem}[商环]
$(R/I, +, \cdot)$ 构成环，称为 $R$ 模 $I$ 的\keyword{商环}。
\end{theorem}

\begin{proof}
加法部分的全部性质由商群理论保证；乘法只需验证\keyword{良好定义性}。
设 $a + I = a' + I$，$b + I = b' + I$，则 $a' = a + i$，$b' = b + j$
（$i,j\in I$），于是
\[
  a'b' = (a+i)(b+j) = ab + aj + ib + ij.
\]
由吸收性，$aj, ib, ij\in I$，故 $aj + ib + ij\in I$，从而
\[
  a'b' + I = ab + I,
\]
即乘积与代表元选取无关。
\end{proof}

% ===========================================================================
\section{同构定理}\label{sec:iso-theorems}

\subsection{第一同构定理（同态基本定理）}

\begin{theorem}[环的第一同构定理]
设 $f\colon R\to R'$ 为环同态，则 $f$ 可分解为满同态、同构、单同态的
复合：
\[
  R \xrightarrow{\ \pi\ } R/\ker f \xrightarrow{\ \bar f\ }
  \Img f \xrightarrow{\ \iota\ } R',
\]
其中 $\pi\colon a\mapsto a+\ker f$ 是典范投影（满同态），
$\iota\colon\Img f\hookrightarrow R'$ 是包含映射（单同态），
$\bar f\colon a+\ker f\mapsto f(a)$ 是诱导同构。特别地，
\[
  R/\ker f \;\cong\; \Img f.
\]
\end{theorem}

\begin{proof}[证明要点]
$\bar f$ 的良好定义性：若 $a' = a + i$（$i\in\ker f$），则
\[
  f(a') = f(a) + f(i) = f(a) + 0 = f(a).
\]
$\bar f$ 显然满；又 $\bar f(a+\ker f) = 0 \iff f(a) = 0 \iff a\in\ker f$，
故单，从而是同构。
\end{proof}

\subsection{理想的运算}

\begin{proposition}[理想之交仍为理想]
任意多个理想的交仍是理想（由定义逐条验证即得）。
\end{proposition}

\begin{definition}[生成理想与主理想（Principal Ideal）]
设 $S\subseteq R$，所有包含 $S$ 的理想之交是包含 $S$ 的\keyword{最小
理想}，称为 \keyword{$S$ 生成的理想}，记 $\langle S\rangle$。由单个
元素生成的理想 $\langle a\rangle$ 称为\keyword{主理想}。
\end{definition}

\begin{proposition}[生成理想的显式刻画]
\[
  \langle S\rangle = \left\{ \sum_{i=1}^{n} x_i\, s_i\, y_i
    \ \middle|\ n\in\N,\ x_i, y_i\in R,\ s_i\in S \right\}.
\]
特别地，$\langle a\rangle = \left\{\sum_i x_i\, a\, y_i\right\}$。
若 $R$ 交换（含幺），则
\[
  \langle a\rangle = aR = \{\,ar : r\in R\,\}.
\]
\end{proposition}

\begin{proof}
右端 $\subseteq \langle S\rangle$：由吸收性每项
$x_i s_i y_i\in\langle S\rangle$，有限和亦然。右端
$\supseteq \langle S\rangle$：右端对加减法封闭，且被任意环元素左乘
（右乘）后形状不变（乘进去的元素吸收进 $x_i$ 或 $y_i$），故右端本身是
包含 $S$ 的理想。
\end{proof}

\begin{definition}[理想的和]
设 $I, J$ 为理想，定义
\[
  I + J = \{\,x + y \mid x\in I,\ y\in J\,\}.
\]
\end{definition}

\begin{proposition}
$I + J$ 仍是理想。
\end{proposition}

\begin{proof}
加减封闭：$(x_1 + y_1)\pm(x_2 + y_2) = (x_1\pm x_2) + (y_1\pm y_2)\in I+J$。
吸收性：对任意 $r\in R$，
\[
  r(x+y) = rx + ry, \qquad rx\in I,\ ry\in J,
\]
故 $r(x+y)\in I+J$；另一侧同理。
\end{proof}

\begin{definition}[理想的积]
字面乘积集合 $\{xy \mid x\in I, y\in J\}$ 一般不是理想，定义 $IJ$ 为它
\keyword{生成的理想}。
\end{definition}

\begin{example}[$\Z$ 的理想都是主理想]
设 $I$ 是 $\Z$ 的理想，令 $d = \gcd\{\,a \mid a\in I\,\}$。辗转相减始终
在 $I$ 内部进行，故 $d\in I$，从而 $(d)\subseteq I$；另一方面 $I$ 中
元素都是 $d$ 的倍数，即 $I\subseteq(d)$。所以 $I = (d)$：$\Z$ 的每个
理想都是主理想。由此
\[
  (a) + (b) = \big(\gcd(a,b)\big), \qquad
  (a)\cap(b) = \big(\mathrm{lcm}(a,b)\big), \qquad
  (a)(b) = (ab).
\]
这些好性质只对整数环成立，一般的环没有这么好的事情。
\end{example}

\subsection{同构定理（续）}

\begin{definition}[同构]
既是同态又是双射的映射称为\keyword{同构}；同构的两个对象在代数意义上
等价。
\end{definition}

\block{商同态}
对任意理想 $I$，自然满同态
\[
  \pi\colon R\longrightarrow R/I, \qquad x\mapsto x + I,
\]
的核恰为 $I$。

\begin{theorem}[第三同构定理]
设 $I\subseteq J$ 是 $R$ 的理想，则
\[
  (R/I)\,\big/\,(J/I) \;\cong\; R/J.
\]
\end{theorem}

\begin{proof}
构造 $\varphi\colon R/I\to R/J$，$x + I\mapsto x + J$（由
$I\subseteq J$ 知良定义）。它是满射，且
\[
  \ker\varphi = \{\,x + I : x\in J\,\} = J/I,
\]
由第一同构定理即得。
\end{proof}

\begin{theorem}[对应定理]
设 $\varphi\colon R\to R'$ 是满同态，则有双射
\[
  \{\,R\text{ 中包含 }\ker\varphi\text{ 的理想}\,\}
  \;\longleftrightarrow\;
  \{\,R'\text{ 的理想}\,\},
\]
正向取像 $I\mapsto\varphi(I)$，反向取原像
$J'\mapsto\varphi^{-1}(J')$。特别地（取 $\varphi$ 为商同态），$R/I$
的理想与 $R$ 中包含 $I$ 的理想一一对应。
\end{theorem}

% ===========================================================================
\section{极大理想与素理想}\label{sec:maximal-prime}

\subsection{单环与域的刻画}

\begin{definition}[平凡理想、单环]
零理想 $\{0\}$ 与整个环 $R$ 称为\keyword{平凡理想}；没有非平凡理想的环
称为\keyword{单环}（类比单群：没有非平凡正规子群的群）。
\end{definition}

\begin{theorem}[含幺交换环是域的刻画]
设 $R$ 为含幺交换环，则
\[
  R\text{ 是域} \iff R\text{ 没有非平凡理想（}R\text{ 是单环）}.
\]
\end{theorem}

\begin{proof}
（$\Rightarrow$）域的理想只要包含非零元 $a$，由 $a$ 可逆，
$1 = a^{-1}a$ 也落入该理想，于是理想等于整个环。

（$\Leftarrow$）任取 $a\in R$，$a\ne 0$。因 $R$ 交换，
$(a) = \{\,ra : r\in R\,\}$。$(a)\ne\{0\}$，而理想只有平凡的两个，故
$(a) = R\ni 1$，即存在 $r$ 使 $ra = 1$，$a$ 可逆。每个非零元可逆，
故 $R$ 是域。
\end{proof}

\subsection{极大理想（Maximal Ideal）}

\begin{definition}[极大理想]
理想 $M\subsetneq R$ 称为\keyword{极大理想}，若任何包含 $M$ 的理想
$I$ 只有 $I = M$ 或 $I = R$ 两种可能。
\end{definition}

\begin{theorem}[极大理想的刻画]
设 $R$ 为含幺交换环，$M\subsetneq R$ 为理想，则
\[
  M\text{ 是极大理想} \iff R/M\text{ 是域}.
\]
\end{theorem}

\begin{proof}
由对应定理，$R/M$ 的理想与 $R$ 中包含 $M$ 的理想一一对应。

（$\Rightarrow$）$M$ 极大，故包含 $M$ 的理想只有 $M$ 与 $R$ 两个，
从而 $R/M$ 没有非平凡理想。又 $R/M$ 交换，由上面的定理，$R/M$ 是域。

（$\Leftarrow$）$R/M$ 是域，由上面的定理它没有非平凡理想，故 $R$ 中
包含 $M$ 的理想只有 $M$ 与 $R$ 两个——这正是 $M$ 极大的定义。
\end{proof}

\subsection{素理想（Prime Ideal）}

\begin{definition}[素理想]
理想 $P\subsetneq R$ 称为\keyword{素理想}，若对任意 $a,b\in R$，
\[
  ab\in P \implies a\in P\ \text{或}\ b\in P.
\]
\end{definition}

\begin{example}
在 $\Z$ 中，素数 $p$ 的倍数构成素理想 $(p)$：若 $ab$ 是 $p$ 的倍数，
则 $a, b$ 中至少一个是 $p$ 的倍数。
\end{example}

\begin{theorem}[素理想的刻画]
设 $R$ 为含幺交换环，$P\subsetneq R$ 为理想，则
\[
  P\text{ 是素理想} \iff R/P\text{ 是整环}.
\]
\end{theorem}

\begin{proof}
商环 $R/P$ 中（零元即 $P$ 本身）：
\[
  (a + P)(b + P) = 0 \iff ab + P = P \iff ab\in P, \qquad
  a + P = 0 \iff a\in P.
\]

（$\Rightarrow$）若 $P$ 是素理想：$(a+P)(b+P) = 0$ 即 $ab\in P$，
由素理想定义 $a\in P$ 或 $b\in P$，即 $a+P = 0$ 或 $b+P = 0$。故 $R/P$
无零因子，是整环。

（$\Leftarrow$）若 $R/P$ 是整环：$ab\in P$ 即 $(a+P)(b+P) = 0$，
由无零因子得 $a+P = 0$ 或 $b+P = 0$，即 $a\in P$ 或 $b\in P$。故 $P$
是素理想。
\end{proof}

\begin{remark}
两个方向的证明其实是同一串等价的正读与倒读：$ab\in P$ 与
“$(a+P)(b+P)=0$”本就是一回事，素理想的定义正是“商环无零因子”的
逐元素写法。
\end{remark}

% ===========================================================================
\section{中国剩余定理与分式域}\label{sec:crt-frac}

\subsection{中国剩余定理（Chinese Remainder Theorem）}

\begin{definition}[互素理想]
理想 $I, J$ 称为\keyword{互素}，若 $I + J = R$（等价地，$1\in I + J$）。
\end{definition}

回顾整数版本：$(a) + (b) = (\gcd(a,b))$，故 $(a)$ 与 $(b)$ 互素当且仅当
$\gcd(a,b) = 1$。

\begin{proposition}
设 $R$ 为含幺交换环，$I, J$ 互素，则
\[
  I\cap J = IJ.
\]
\end{proposition}

\begin{proof}
（$\supseteq$）$IJ$ 中元素 $xy$（$x\in I$, $y\in J$）既属于 $I$ 又属于
$J$，故 $IJ\subseteq I\cap J$。

（$\subseteq$）由互素，写成 $1 = i + j$（$i\in I$, $j\in J$）。任取
$a\in I\cap J$，则
\[
  a = a\cdot 1 = ai + aj,
\]
其中 $ai\in IJ$（$i\in I$ 且 $a\in J$），$aj\in IJ$（$a\in I$ 且
$j\in J$），故 $a\in IJ$。
\end{proof}

\begin{example}[无幺元时命题不成立]
取 $R = 2\Z$（无幺元），$I = J = R$，则 $I\cap J = 2\Z$ 而 $IJ = 4\Z$，
两者不等。
\end{example}

\begin{theorem}[中国剩余定理（CRT）]
设 $R$ 为含幺交换环，$I, J$ 互素，则
\[
  R/(I\cap J) \;\cong\; R/I\times R/J
\]
（由上一命题，左端分母也可写成 $IJ$）。
\end{theorem}

\begin{proof}
构造环同态
\[
  \varphi\colon R\longrightarrow R/I\times R/J, \qquad
  x\mapsto (x + I,\ x + J).
\]

\keyword{核：}$\varphi(x) = 0 \iff x\in I$ 且 $x\in J$，故
$\ker\varphi = I\cap J$。

\keyword{满：}任取 $(a + I,\ b + J)$。写 $1 = i + j$（$i\in I$,
$j\in J$），取 $x = aj + bi$。则
\[
  x + I = aj + I = a(1-i) + I = a + I \quad(\text{因 }ai\in I),
\]
同理 $x + J = b + J$。故 $\varphi(x) = (a+I,\ b+J)$。

由第一同构定理，$R/(I\cap J)\cong R/I\times R/J$。
\end{proof}

\begin{remark}[整数版本]
$p, q$ 互素时（取 $I = p\Z$, $J = q\Z$），
\[
  \Z/pq\Z \;\cong\; \Z/p\Z\times\Z/q\Z.
\]
\end{remark}

\subsection{单位、相伴与单位群（Group of Units）}

\begin{definition}[单位]
元素 $u\in R$ 称为\keyword{单位}（可逆元），若存在 $v\in R$ 使
$uv = vu = 1$。注意一般的环中，除 $0$ 之外还可能有其他元素也不是单位
（这与域不同）。
\end{definition}

\begin{definition}[相伴关系]
在交换环中，定义 $r\sim s \iff$ 存在单位 $u$ 使 $r = us$。这是一个
等价关系。
\end{definition}

\begin{definition}[单位群（Group of Units）]
$R$ 的全体单位
\[
  R^\times = \{\,u\in R \mid u\text{ 是单位}\,\}
\]
在乘法下构成群，称为 $R$ 的\keyword{单位群}。
\end{definition}

\begin{example}
$\Z^\times = \{\pm 1\}$；域 $k$ 上多项式环的单位恰为非零常数多项式：
$k[x]^\times = k\setminus\{0\}$（次数 $\ge 1$ 的多项式均不可逆）。
\end{example}

\subsection{整环的分式域（Field of Fractions）}

\block{分式域的构造}
设 $R$ 为整环。形式上考虑“分数”
\[
  \frac{a}{b}, \qquad a\in R,\ b\in R\setminus\{0\}.
\]
分数的表示不唯一（$\frac{a}{b}$ 与 $\frac{2a}{2b}$ 是同一个分数），
模去等价关系：
\[
  \frac{a}{b}\sim\frac{c}{d} \iff ad = bc
\]
（交叉相乘相等）。等价类全体记为 $\Frac(R)$，按普通分数的法则定义
加法与乘法。

\begin{theorem}[分式域]
设 $R$ 为整环，则 $\Frac(R)$ 是域，称为 $R$ 的\keyword{分式域}。
\end{theorem}

\begin{proof}[证明要点]
关键是每个非零元有逆：
\[
  \frac{a}{b}\ne 0 \iff a\ne 0,\qquad
  \text{此时其逆为}\ \frac{b}{a};
\]
乘法幺元为 $\frac{1}{1}$（的等价类）。
\end{proof}

\begin{example}
\[
  \Frac(\Z) = \Q; \qquad
  \Frac\big(k[x]\big) = k(x)\ (\text{有理函数域}).
\]
\end{example}

\subsection{特征（Characteristic）与有限域 \texorpdfstring{$\F_p$}{Fp}}

\begin{definition}[特征（Characteristic）]
环 $R$ 的\keyword{特征} $\charr R$ 是使
\[
  \underbrace{1 + 1 + \cdots + 1}_{n\text{ 次}} = 0
\]
首次成立的最小正整数 $n$；若 $1$ 无论加多少次都不等于 $0$，则称特征为
$0$。等价地，凡满足 $m\cdot 1 = 0$ 的整数 $m$ 都被 $\charr R$ 整除。
\end{definition}

\begin{proposition}
整环（特别地，域）的特征若非零，必为素数。
\end{proposition}

\begin{example}[有限域 $\F_p$]
对任意素数 $p$，模 $p$ 剩余类构成域
\[
  \F_p = \Z/p\Z,
\]
它是最小的特征为 $p$ 的域。
\end{example}

\block{构造（无穷的特征 $p$ 域）}
从有限域 $\F_p$ 出发，先取多项式环 $\F_p[x]$（这是一个整环），再取其
分式域：
\[
  \F_p(x) = \Frac\big(\F_p[x]\big),
\]
便得到一个元素个数无穷、特征为 $p$ 的域。


\end{document}
